\documentclass[11pt]{article}
\usepackage[margin=1in]{geometry}
\usepackage{amsmath}
\usepackage{booktabs}
\usepackage{hyperref}

\title{A Quantum-Inspired Energy-Based Scheduling Framework for Hardware-Aware NPU Optimization}
\author{Yagnesh Kumar Koduru}
\date{2026}

\begin{document}
\maketitle

\section{Introduction}
The project started from a practical issue we kept seeing in NPU scheduling: a schedule that looks fine at graph level can still perform poorly once SRAM pressure, bandwidth bursts, and queue backlogs are accounted for. A purely local heuristic may serialize work too early, evict data that should have stayed on-chip, or create memory traffic spikes that later show up as stalls.

So instead of building a single ``smart'' scheduler, we built a layered system. First, multiple classical strategies produce candidate topological orders. Then a hardware-aware cost model scores them. Finally, an adaptive penalty mechanism (APR) reshapes constraint importance over rounds, and a quantum-style refinement stage searches harder neighborhoods from a warm start.

\section{What Is Actually Implemented}
The repository is modular, and each file has a clear role:
\begin{itemize}
    \item \texttt{memory\_hierarchy.py}: SRAM residency simulation with eviction policy, prefetch credit, spill tracking, bank conflict cycles, and capacity violations.
    \item \texttt{bandwidth\_estimator.py}: read/write channel utilization, burst-window averaging, backlog carry-over, congestion and pipeline stall estimation.
    \item \texttt{fusion\_logic.py}: rule-based pair/triplet fusion gain estimation (e.g., conv--batchnorm--relu style chains).
    \item \texttt{cost\_model.py}: combines memory, bandwidth, fusion, and parallelism terms into total cost plus penalty terms.
    \item \texttt{energy\_model.py}: energy objective builder with unary cost plus pairwise interaction terms (reuse/fusion rewards, bandwidth and memory conflict penalties).
    \item \texttt{penalty\_tuner.py}: APR update loop using violation rate, cost impact, and recurrence frequency.
    \item \texttt{scheduling\_engine.py}: Greedy, Lookahead, Beam Search, and Simulated Annealing order generation.
    \item \texttt{run\_experiment.py}: end-to-end orchestration for both legacy cost-based and new energy-based formulations, including X-factor evaluation and comparison export.
\end{itemize}

\section{Cost Model (Code-Aligned)}
The objective in \texttt{cost\_model.py} is not a single scalar heuristic; it is a structured sum of physically motivated terms:
\[
\begin{aligned}
C_{\text{raw}} =
&\;w_d\big(\text{dram\_access} + 0.08\cdot\text{spill\_count}\cdot\max(1,\text{avg\_sram\_util})\big) \\
&+ w_r\big(\text{sram\_reuse\_loss} + 0.45\cdot\text{bank\_conflict\_cycles}\big) \\
&+ w_b\big(\text{bandwidth\_congestion} + 0.35\cdot\text{backlog\_pressure}\big) \\
&+ w_s\big(\text{pipeline\_stalls} + 0.22\cdot\text{bank\_conflict\_cycles}\big) \\
&- w_f\cdot\text{fusion\_gain} + w_p\cdot\text{parallelism\_loss}
\end{aligned}
\]

Then penalties are added:
\[
\begin{aligned}
C_{\text{penalty}} =
&\;140\,M_{\text{sram}}v_{\text{sram}} + 100\,M_{\text{bw}}v_{\text{bw}} + 900\,M_{\text{dep}}v_{\text{dep}} \\
&+ 70\,M_{\text{dram}}\max(0, v_{\text{dram}}-1)
+ 90\,M_{\text{bank}}v_{\text{bank}}
+ 65\,M_{\text{imb}}v_{\text{imb}}
\end{aligned}
\]
\[
C_{\text{total}} = C_{\text{raw}} + C_{\text{penalty}}
\]

\textbf{Why these terms were chosen:}
\begin{itemize}
    \item \textbf{DRAM term}: captures expensive off-chip movement and spill side effects.
    \item \textbf{SRAM reuse loss}: captures re-fetch pain caused by eviction or missed reuse windows.
    \item \textbf{Bandwidth congestion + stalls}: models burst pressure and how it directly delays execution.
    \item \textbf{Fusion gain}: rewards schedules that place fusible operator chains adjacently.
    \item \textbf{Parallelism loss}: discourages orders that serialize the frontier even when parallel-ready work exists.
\end{itemize}

In this implementation, \texttt{parallelism\_loss} is critical-path-aware: when many nodes are ready but one branch is chosen repeatedly, the model adds serialization and slack penalties so the scheduler does not accidentally collapse available concurrency.

\section{Energy-Based Scheduling Formulation}
The scheduling objective is reformulated as an energy minimization problem so that local operator interactions are optimized jointly with the global hardware-aware unary cost:
\[
E_{\text{total}} = E_{\text{unary}} + E_{\text{pairwise}}
\]
with
\[
E_{\text{unary}} = C_{\text{total}}
\]
and
\[
\begin{aligned}
E_{\text{pairwise}} =
&\;\lambda_{\text{bw}}P_{\text{bandwidth\_spike}}
+ \lambda_{\text{mem}}P_{\text{memory\_conflict}} \\
&- \lambda_{\text{reuse}}R_{\text{data\_reuse}}
- \lambda_{\text{fusion}}R_{\text{fusion\_opportunity}}.
\end{aligned}
\]

\textbf{Why pairwise interactions matter:}
\begin{itemize}
    \item Two schedules can have similar unary cost but very different local ordering quality.
    \item Adjacent operator pairs determine immediate data reuse windows; bad adjacency forces extra movement and evictions.
    \item Fusion opportunities are inherently pair/triplet-local and cannot be fully captured by unary node terms.
    \item Bandwidth spikes and bank conflicts are often burst effects caused by consecutive placements, not by one node alone.
\end{itemize}

\textbf{How this differs from traditional scheduling:}
\begin{itemize}
    \item Traditional formulations optimize mostly node-level or aggregate costs.
    \item The new formulation explicitly models neighbor interactions as first-class terms in the objective.
    \item Both classical search and QAOA-style refinement now optimize the same energy objective, enabling quantum-inspired local/global balancing.
\end{itemize}

\section{APR in This Project}
APR is implemented in \texttt{penalty\_tuner.py}. Each round updates penalty multipliers from three signals:
\begin{itemize}
    \item \texttt{violation\_rate}: how much a constraint is being violated right now.
    \item \texttt{cost\_impact}: how strongly that constraint's region contributes to total cost.
    \item \texttt{constraint\_frequency}: how often the same constraint keeps appearing across rounds.
\end{itemize}

The update is:
\[
\text{growth} = \exp(\text{violation\_rate} + \text{cost\_impact} + \text{constraint\_frequency})
\]
with practical guards:
\begin{itemize}
    \item exponential growth is clipped by \texttt{max\_growth = 2.4}
    \item momentum blending uses \texttt{0.18} to reduce oscillation
    \item penalties are clipped to \([0.2, 14.0]\)
    \item both violation and cost signals are EMA-smoothed (decay \texttt{0.72})
\end{itemize}

\textbf{Why this matters:} violation alone is too reactive, and cost impact alone can ignore rare but dangerous constraints. Frequency keeps the system from forgetting recurring bottlenecks.

\section{Memory and Bandwidth Modeling Choices}
\subsection*{Memory hierarchy}
The SRAM simulator keeps an ordered tensor cache and evicts using a score that balances remaining consumers, tensor size, and age. It also tracks prefetch credits, bank conflicts, peak SRAM usage, spill count, and idle-cycle impact from evictions.

This design gives the scheduler feedback on \emph{why} a schedule is bad: not just ``DRAM is high,'' but whether the issue comes from poor reuse timing, bank pressure, or outright capacity overrun.

\subsection*{Bandwidth estimator}
Bandwidth uses separate read/write capacities with burst windows and decayed backlog. Overflow contributes to both congestion and pipeline stalls, and read-write asymmetry amplifies penalties.

Without this, the optimizer can choose orders that look cheap in memory terms but create traffic waves that stall compute later.

\section{Real Experiment Flow (\texttt{run\_experiment.py})}
The script follows this pipeline:
\begin{enumerate}
    \item Load \texttt{config.yaml} and \texttt{example\_workload.json} (31-node DAG).
    \item Instantiate memory, bandwidth, fusion, cost-model, and energy-model modules.
    \item Run two full suites: legacy cost-based objective and new energy-based objective. Each suite executes Greedy, Lookahead, Beam Search, Simulated Annealing, Quantum (QAOA-style), and Quantum+APR.
    \item For each strategy, run multiple seeded trials and pick the best via
    \[
    \text{rank\_score} = \phi(\text{schedule}) - 120 \cdot \text{feasibility},
    \]
    where \(\phi\) is \(\text{total\_cost}\) for legacy runs and \(\text{total\_energy}\) for energy-based runs.
    \item Run Quantum (QAOA-style) refinement starting from the annealing schedule.
    \item Warm APR using strong classical candidates (lookahead/beam/anneal and APR-warm SA), then choose an APR seed order by feasibility-first ranking.
    \item Run Quantum+APR for multiple rounds, update penalties each round, and select the best trial.
    \item Save \texttt{outputs/schedules.json}, \texttt{outputs/metrics.txt}, and \texttt{outputs/explanations.txt} with old-vs-new deltas and X-factor summaries.
\end{enumerate}

\section{Main Results and Interpretation}
Latest run from the current code (\texttt{python run\_experiment.py}):

\begin{table}[h]
\centering
\begin{tabular}{lrrrr}
\toprule
Track & Old Cost & New Cost & Cost Delta & Feasible Delta (\%) \\
\midrule
Lookahead & 4168.69 & 4170.84 & -2.15 & +0.00 \\
Quantum (QAOA) & 4439.67 & 4441.83 & -2.16 & +0.00 \\
Quantum + APR & 4168.69 & 4216.92 & -48.23 & +3.23 \\
\bottomrule
\end{tabular}
\caption{Old-vs-new formulation comparison (cost and feasibility).}
\end{table}

\begin{table}[h]
\centering
\begin{tabular}{lrrr}
\toprule
X-factor Track & Baseline Cost & Optimized Cost & X \\
\midrule
Classical methods & 5669.65 & 4168.69 & 0.2647 \\
Quantum & 5669.65 & 4439.67 & 0.2169 \\
Quantum + new formulation & 5669.65 & 4216.92 & 0.2562 \\
\bottomrule
\end{tabular}
\caption{Efficiency factor \(X = (\text{baseline\_cost} - \text{optimized\_cost})/\text{baseline\_cost}\).}
\end{table}

\textbf{Interpretation instead of just ranking:}
\begin{itemize}
    \item Cost deltas between old and new formulations are small for Lookahead and Quantum baseline (about 0.05\%).
    \item The energy-based Quantum+APR run improves feasibility by 3.23 points and slightly improves latency, showing a stronger constraint profile.
    \item X-factor remains strongly positive across classical, quantum, and quantum+new tracks versus the legacy baseline cost.
\end{itemize}

\section{Ablation Study}
I ran targeted ablations to understand what each subsystem contributes.

\begin{table}[h]
\centering
\begin{tabular}{lrrr}
\toprule
Ablation Setting & Cost & Latency & Feasible(\%) \\
\midrule
Baseline (Lookahead, full model) & 4168.69 & 3390.81 & 54.83 \\
Remove APR (use Quantum without APR) & 4439.67 & 3443.75 & 54.83 \\
Remove fusion gain (same Lookahead order, fusion off) & 4414.48 & 3390.81 & 54.83 \\
Remove bandwidth model from search objective$^\dagger$ & 4469.35 & 3496.34 & 51.61 \\
\bottomrule
\end{tabular}
\caption{Ablation effects measured on current code.}
\end{table}

$^\dagger$ For this ablation, scheduling was done with bandwidth terms/penalties disabled, then evaluated using the full baseline model.

\textbf{What this tells us:}
\begin{itemize}
    \item \textbf{Without APR:} total cost is lower than Quantum+APR in this run, but APR still improves latency and traffic smoothness. Here APR is conservative and over-penalizes.
    \item \textbf{Without fusion:} cost increases by 245.79 (about 5.9\%), showing fusion offsets a meaningful portion of memory+compute overhead.
    \item \textbf{Without bandwidth modeling:} cost increases by about 7.2\%, latency by about 3.1\%, and feasibility drops 3.22 points. So bandwidth awareness is not optional.
\end{itemize}

\section{Design Insights}
\begin{itemize}
    \item Using multiple classical search styles is better than trusting one heuristic. Different methods fail in different ways.
    \item Separating memory and bandwidth models gives actionable diagnostics; combined black-box penalties were much harder to tune.
    \item The feasibility-aware trial rank (cost minus feasibility bonus) is simple but useful for choosing stable candidates.
    \item APR needs clipping/momentum. Without these guards, penalty multipliers can oscillate or explode.
    \item Quantum-style refinement works best as a local improver on strong warm starts, not as a standalone scheduler.
\end{itemize}

\section{Limitations}
\begin{itemize}
    \item Small-scale simulation only (single synthetic 31-node workload in current setup).
    \item No real NPU hardware-in-the-loop validation.
    \item Memory and bandwidth models are detailed but still approximate, not cycle-accurate microarchitecture simulation.
    \item Quantum stage is simulator-style, so claims are about search behavior, not hardware quantum speedup.
\end{itemize}

\section{Conclusion}
This project is strongest as a systems-oriented scheduling study: a real pipeline, explicit hardware signals, and measurable trade-offs. The main takeaway from the current code is that hardware-aware modeling and search diversity provide consistent gains, while APR gives useful constraint control but still needs better calibration to avoid over-penalization on this workload.

\end{document}
